\documentclass[11pt]{amsart}
\usepackage[T1]{fontenc}
\usepackage[utf8]{inputenc}
\usepackage{amssymb}
\usepackage[margin=1.05in]{geometry}
\usepackage[hidelinks]{hyperref}
\newcommand{\stamptime}{\the\numexpr\time/60\relax:\ifnum\numexpr\time-60*(\time/60)\relax<10 0\fi\the\numexpr\time-60*(\time/60)\relax}
\AddToHook{shipout/foreground}{\put(0.5\paperwidth,-\paperheight+0.55cm){\makebox[0pt]{\tiny\ttfamily unit/doc/\jobname.tex \ --- \ compiled \today, \stamptime}}}
\newtheorem{theorem}{Theorem}
\newtheorem{proposition}[theorem]{Proposition}
\newtheorem{lemma}[theorem]{Lemma}
\newtheorem{corollary}[theorem]{Corollary}
\theoremstyle{definition}
\newtheorem{definition}[theorem]{Definition}
\newtheorem{axiom}[theorem]{Axiom}
\theoremstyle{remark}
\newtheorem{remark}[theorem]{Remark}
\newcommand{\Om}{\Omega}
\newcommand{\Map}{\mathrm{Map}_*}
\newcommand{\Z}{\mathbb{Z}}
\newcommand{\code}[1]{\texttt{#1}}
\newcommand{\cT}{\mathsf{isTrunc}}
\newcommand{\cC}{\mathsf{isConn}}
\newcommand{\refl}{\mathsf{refl}}
\newcommand{\ap}{\mathsf{ap}}
\newcommand{\lp}{\Omega}
\newcommand{\Hom}{\mathrm{Hom}}
\newcommand{\hc}[2]{[#1,#2]_*}
\newcommand{\Fib}{\mathrm{Fib}}
\newcommand{\fib}{\mathrm{fib}}
\newcommand{\tr}[1]{\|#1\|}
\newcommand{\cls}[1]{[#1]}
\newcommand{\kb}{\bar\kappa}
\newcommand{\Fh}{\hat F}
\newcommand{\sq}{\mathsf{sq}}
\newcommand{\cE}{\mathcal E}
\newcommand{\Unl}{\mathsf{Unloop}}
\newcommand{\cj}{\mathrm{conj}}
\newcommand{\ucj}{\mathrm{unconj}}
\title{The marked universal property of $BR$, and $\pi_4(S^3)=\Z/2$}
\date{28 October 2026}
\begin{document}
\maketitle

\noindent\textbf{Purpose and status.} Step \S2e of
\code{doc/pi3-s2-postnikov-handoff-10.txt} (step~3 of handoff~9). We
specify a class $s\in\hc{S^3}{BR}$ and prove that
\[s^*:\hc{BR}Y\to\hc{S^3}Y\]
is a bijection for every pointed 2-connected 4-type $Y$
(Theorem~\ref{t:up}); the two instances $Y=BR$ and $Y=\tr{S^3}_4$ then give
$BR\simeq\tr{S^3}_4$ on classes and $\pi_4(S^3)\cong\pi_3(R)\cong\Z/2$ as
groups (Theorem~\ref{t:final}). \emph{Draft, written before any code. Read by the user (28 Oct), corrections incorporated:
Lemma~\ref{l:Epre} made explicit along \code{ENat.e-nat} with Lemma~\ref{l:ccr};
Corollary~\ref{c:conv} (Lemma~\ref{l:N} settles the numerical converse); the
diagonal evaluation in Proposition~\ref{p:Q}. \emph{Refereed} independently
(a fresh reader given the note and the library, not the conclusion; 28 Oct): sound
with repairs, none fatal; incorporated: a generic two-sided \code{Conj} (the
single-degree \code{Climb.conj-unconj} does not apply), the lock bridges, the index of
\code{ConnPi1} in Proposition~\ref{p:trunc4}, the direction of \code{loop-homotopy}, the
obligations for $\pi_4$ on classes, \code{Triple3}, the input of the last step of
\code{Rec4}, \code{unique-loop1}.} Numbered statements are meant to be declarations
(\S\ref{s:corr} proposes names).

\smallskip\noindent\textbf{Assumption.} The \emph{$S^3$ axiom}
(Axiom~\ref{a:s3}): the universal property of $S^3$, stated as a segment
parameter telescope as \code{SphereAxiom} does for $S^2$, \emph{not} as a
Theory parameter. It is the only assumption beyond the foundations of the
development (univalence, images; function extensionality derived). The
$S^2$ axiom is \emph{not} used: the generator of $\pi_3K(\Z,3)$ is built from the
circle property (Definition~\ref{d:sigma3}), and everything about $R$, $BR$,
$\kappa$ is axiom-free (\code{doc/cross-term.pdf}, \code{H4Proof.h4},
\code{FibreH3}).

\smallskip\noindent\textbf{The decisive point.} Constructing $BR$
(\code{DeloopR}) proves that a delooping of $R=\Fib\kb$ \emph{exists}; it says
nothing about maps \emph{out of} $BR$. A map $BR\to Y$ into a target presented
as $Y\simeq\Fib K$, $K:K(A,3)\to_*K(B,5)$, exists over a given map to
$K(A,3)$ exactly when an obstruction vanishes, and the obstruction is a class
$[H]=\Phi_B(h)\in\hc{K_2}{K(B,4)}\cong B$. The fibre sequence of $BR$ kills
it only if $h$ factors through $\rho:\Z\to\Z/2$, i.e.\ if $2h=0$. So
the proof must show: \emph{for every target, the obstruction is 2-torsion.}
The cross-term proof does not give this: it shows that \emph{our} $\kb$ has a
vanishing cross term because $2=0$ in $\Z/2$; the general statement ``if
$e_2;F;e_4^{-1}$ deloops, then $2g=0$'' is the converse of
\code{doc/cross-term.pdf}, Proposition~21 (the handoff calls it 20), which is open there and is
\emph{not} used as established. We prove it here, as Lemma~\ref{l:N}
(``necessity'') and Corollary~\ref{c:conv}, by a route that avoids the obstacle of
\S7 of that note (normalising a cell); since $\Unl f=\sum_{\hat G}(f=\lp\hat G)$, an
inhabitant already supplies the map and the equality Lemma~\ref{l:N} needs:
\begin{itemize}
\item the obstruction $H$ is the looping of a map $K(\Z,3)\to_*K(B,5)$ \emph{by
  construction} (it is $\ucj(\lp M)$ for $M=G;K$), so no cell of
  \code{Unloop} has to be normalised;
\item $\lp M$ commutes with squaring of loops (\code{loop-mult}), and squaring
  is $K(2,2)$ on the source and $K(2_B,4)$ on the target (\code{Doubling}), so
  $K(2,2);H$ and $H;K(2_B,4)$ are pointwise equal; this is a statement about
  underlying maps into $K(B,4)$, where pointings are free (merely);
\item a computation with the bipointed classification (Theorem~B of
  \code{doc/h4.pdf}) gives $\cls{K(2,2);\kappa}=\cls{\kappa;K(4,4)}$
  (Proposition~\ref{p:Q}: $(2x)^2=4x^2$), hence $\Phi_B(4h)=\Phi_B(2h)$ and
  $2h=0$ by injectivity of $\Phi_B$.
\end{itemize}
The only new ingredient of the computation is the naturality of the invariant
$\cE$ under precomposition in the \emph{second} variable
(Lemma~\ref{l:Epre}); first-variable doubling is \code{ETwo.e-two}.

\smallskip\noindent\textbf{Why at the level of $BR$, not of $R$.} One could
loop everything ($\hc{BR}Y\cong\hc R{\Om Y}$, Lemma~\ref{l:loopeq3}) and run
\code{MapFib} with source $R$. But then the target is only a 1-connected
3-type $\Om Y$ \emph{with the property} of being a loop space, and the
obstruction would have to be extracted from an equivalence $\Om Y\simeq\Fib k$,
which is the open converse in disguise. Presenting the target one level up
(Theorem~\ref{t:pres4}) makes its $k$-invariant a map $K:K(A,3)\to_*K(B,5)$,
so every obstruction is a looping by construction. The proof is therefore
\code{MapFib} with every index raised by one, plus Lemma~\ref{l:N} in the
proof of (E) for the source $BR$ (Proposition~\ref{p:EBR}).

\smallskip\noindent\textbf{Restricted coefficients.} The coefficient map of
$BR$ is $\Psi_B:\Hom(\Z/2,B)\to\hc{K_2}{K(B,4)}$,
$g\mapsto\cls{\kb;K(g,4)}=\Phi_B(\rho;g)$. It is injective
(Proposition~\ref{p:Psi}), which is what \code{FibreH3} needs (so
$\hc R{K(B,3)}=0$, Proposition~\ref{p:UR}); and its image contains every
looping class (Lemmas~\ref{l:N}, \ref{l:factor}), which is what (E) needs.
The unrestricted interface (\code{H4}, \code{FibreH3.FibAt},
\code{MapFib.SourceP}) is not instantiated with $\kb$; the restricted
statements are written out below.

\section{Conventions}

As in \code{doc/mapping-fib.pdf}: \code{isTrunc n} is ``$(n-2)$-truncated'',
\code{isConn n} is ``$(n-2)$-connected''; $K^A_m=K(A,m)$ (\code{EMMaps.KK}),
$K_m=K^\Z_m$, $K'_m=K^{\Z/2}_m$; $e_m:\Om K^A_{m+1}\simeq_*K^A_m$; $F;G$ is
\code{ptcomp F G}; $0$ is \code{ptconst}; $\cls F$ is \code{hcl F}, $a;b$ on
classes is \code{compC}; equality of classes is mere equality of pointed maps
(\code{effective}, \code{relates}). $\cj_mF=e_m;F;e_m^{-1}$ and
$\ucj_m(f)=e_m^{-1};f;e_m$ (\code{FibreH3.Climb.unconj}; $\cj\circ\ucj=\mathrm{id}$,
\code{conj-unconj}). In \S\ref{s:N} and Proposition~\ref{p:EBR}, $\cj$ is used with
\emph{different degrees and groups} on the two sides ($\cj F=e_2;F;e_4^{-1}:\Om K_3\to_*\Om K^B_5$
for $F:K_2\to_*K^B_4$), which the single-degree \code{Lev.conj}, \code{Climb.conj-unconj},
\code{H2.unconj-conj} do not cover. A generic \code{Conj[eA,eB]} for two arbitrary pointed
equivalences $e_A:\Om X\simeq_*A$, $e_B:\Om Y\simeq_*B$ is written first, with
\code{conj-unconj}, \code{unconj-conj} (proofs as \code{Climb.conj-unconj}) and composition
$\cj_{A,C}(F;G)=\cj_{A,B}F;\cj_{B,C}G$ (\code{KFun.conj-comp}, already generic); every
$\cj$ below is an instance;
$\lp K(g,m{+}1)=\cj_mK(g,m)$ (\code{loopK}). $\Phi_B(h)=\cls{\kappa;K(h,4)}$
for $h:\Hom(\Z,B)$ (\code{KZ2Cohomology.Phi}), a bijection (\code{H4Proof.h4}).
For homomorphisms, $g;h$ is \code{hcomp g h} $=h\circ g$; $K(g;h,m)=K(g,m);K(h,m)$
(\code{KFun.Comp.k-comp}). $2_A$, $0$, $\rho$ as in \code{Z2}; $4_\Z=2_\Z;2_\Z$.

Levels: $K^A_m$ is $\cC(m{+}1)$ and $\cT(m{+}2)$. A \emph{pointed 2-connected
4-type} is a pointed $Y$ with $\cC\,4$ and $\cT\,6$. From
\code{doc/cross-term.pdf}: $\kb=\kappa;K(\rho,4)$, $R=\Fib\kb$ (\code{isConn 3},
\code{isTrunc 5}), $\Fh:K_3\to_*K'_5$ with $\gamma:e_2;\kb;e_4^{-1}=\lp\Fh$,
$BR=\Fib\Fh$ pointed by $b_0=(*,\Fh_0)$, $p=\mathrm{fpr}:BR\to_*K_3$,
$\phi:\Om BR\simeq_*R$ (\code{omega-BR}); $BR$ is $\cC\,4$, $\cT\,6$.

\section{The $S^3$ axiom and the 4-truncation of $S^3$}

\begin{axiom}[$S^3$]\label{a:s3}
A type $S^3$ with $\mathsf{base}$ and $\mathsf{surf}:\Om^3(S^3,\mathsf{base})$
such that, for every $X$ of the universe of $S^3$, $\mathrm{ev}_3:(S^3\to
X)\to(x:X)\times\Om^3(X,x)$, $f\mapsto(f\,\mathsf{base},\ap^3_f\mathsf{surf})$
is an equivalence. A segment parameter telescope
\code{(S3, base, surf, ump3)}, as \code{SphereAxiom.Sphere}; not a Theory
parameter.
\end{axiom}

\begin{lemma}[as \code{SphereAxiom}]\label{l:ev3}
$\mathrm{ev}:\Map(S^3,X)\to\Om^3X$, $F\mapsto(\lp^3F)(\mathsf{surf})$, is an
equivalence (\code{pointed-up}: $\Sigma$ over \code{ump3} and contraction of the
singleton, with $\mathrm{ev}(f,\pi)=\pi_*\ap^3_f\mathsf{surf}$ by $J$), and
$\mathrm{ev}(F;G)=(\lp^3G)(\mathrm{ev}F)$ (\code{loop-comp} three times).
\end{lemma}

\begin{lemma}[$S^3$ is 2-connected]\label{l:s3conn}
$S^3$ is $\cC\,4$.
\end{lemma}
\begin{proof}
As \code{SphereTrunc.Simply}, one level up: for $Y$ \code{isTrunc 4},
$\Om^3(Y,y)$ is contractible, so evaluation at $\mathsf{base}$ is an
equivalence $(S^3\to Y)\simeq Y$; hence $V=\tr{S^3}_2$ (\code{TO three}) is
contractible (the identity and the constant map agree after the unit), and
$S^3$ is the fibre of the unit over $|\mathsf{base}|$, which is $\cC\,4$.
\end{proof}

\begin{proposition}[the 4-truncation]\label{p:trunc4}
$T=\tr{S^3}_4$ (\code{TO five}), pointed by $t_0=|\mathsf{base}|$, with
$\eta=(|{-}|,\refl):S^3\to_*T$, satisfies: $T$ is $\cT\,6$ and $\cC\,4$;
$\eta^*:\Map(T,Z)\to\Map(S^3,Z)$ is an equivalence for $Z$ \code{isTrunc 6}
(\code{PtUP.up}); and $\pi_4\eta:\pi_4S^3\to\pi_4T$ is a group isomorphism.
Here $\pi_4(X,x_0)=\pi_3(\Om X,\refl)$ by definition, and $\pi_4F=\pi_3(\lp F)$.
\end{proposition}
\begin{proof}
As \code{SphereTrunc.Theorem15}: connectedness by \code{conn-up} from the
fibres of the unit and Lemma~\ref{l:s3conn}; the fibres of $\lp^k\eta$ are
$\cC(6-k)$ (\code{loop-conn}); \code{ConnPi1} for $\lp^3\eta$ takes the
connectivity of the fibres of its looping $\lp^4\eta$, which are $\cC\,2$
(\code{loop-conn} four times), and makes $\pi_1(\lp^3\eta)=\pi_4\eta$ bijective.
\end{proof}

\section{$K(\Z,3)$ is the 3-truncation of $S^3$}

\begin{lemma}[looping is an equivalence]\label{l:loopeq3}
For $X$ \code{isConn (c+2)} and $Y$ \code{isTrunc (2c+2)}, $\lp:\Map(X,Y)\to
\Map(\Om X,\Om Y)$ is an equivalence (\code{Over.cor11}); used with $c=2$
($X=K_3$ or $BR$, $Y$ a 3-type or 4-type) and $c=1$ ($X=K_2$,
\code{MapFib.SourceS.LoopEq}). For $X$ \code{isConn 4} and $Y$
\code{isTrunc 7}, $\lp$ is an embedding: its fibres $\Unl f$ are propositions
(\code{Levels[2,1].cor10}).
\end{lemma}

\begin{definition}[$\sigma_3$]\label{d:sigma3}
Let $g:\Om K_1$ be the circle generator, $c_3=(\lp^2e_2^{-1})((\lp e_1^{-1})(g)):\Om^3K_3$,
and $\sigma_3=\mathrm{ev}^{-1}(c_3):S^3\to_*K_3$.
\end{definition}

\begin{proposition}\label{p:sigma3}
For every 3-type $(Z,z_0)$ (\code{isTrunc 5}), $\sigma_3^*:\Map(K_3,Z)\to\Map(S^3,Z)$
is an equivalence; hence $\cls{\sigma_3}^*:\hc{K_3}Z\to\hc{S^3}Z$ is a bijection.
\end{proposition}
\begin{proof}
Let $\Psi_3(G)=\mathrm{ev}^\circ\bigl(e_1^{-1};\lp(e_2^{-1};\lp G)\bigr)$, where
$\mathrm{ev}^\circ:\Map(K_1,\Om^2Z)\simeq\Om^3Z$, $N\mapsto(\lp N)(g)$, is the
circle property. $\Psi_3$ is a composite of five equivalences:
$\lp$ (Lemma~\ref{l:loopeq3}, $c=2$; $Z$ is $\cT\,6$ by \code{trunc-suc}),
$e_2^{-1};-$ (\code{pre-pteq}), $\lp$ ($c=1$; $\Om Z$ is $\cT\,4$),
$e_1^{-1};-$, $\mathrm{ev}^\circ$. By Lemma~\ref{l:ev3} and \code{loop-comp},
\[\mathrm{ev}(\sigma_3;G)=(\lp^3G)(c_3)=\bigl(\lp^2(e_2^{-1};\lp G)\bigr)((\lp e_1^{-1})g)
=\bigl(\lp(e_1^{-1};\lp(e_2^{-1};\lp G))\bigr)(g)=\Psi_3(G).\]
So $\mathrm{ev}\circ\sigma_3^*=\Psi_3$ and $\sigma_3^*$ is an equivalence
(\code{cancel-equiv}); classes by \code{pre-equiv}. The $S^2$ axiom is not used.
\end{proof}

\section{Level-3 triples and the lifting lemma}

A \emph{level-3 triple} is $(A,B,K)$ with $K:K^A_3\to_*K^B_5$. Its fibre
$X=\Fib K$, pointed by $(*,K_0)$, with $q=\mathrm{fpr}:X\to_*K^A_3$, is
$\cC\,4$ (fibres $K\,y=*$ are $\cC\,5$; \code{FibConn}) and $\cT\,6$.
$\Om K^B_5$ is $\cC\,5$, $\Om^3K^B_5$ is $\cC\,3$, $\Om^4K^B_5$ is $\cC\,2$,
$\Om X$ is $\cC\,3$; all connected.

\begin{theorem}[lifting, \code{MapFib} Theorem~6 one level up]\label{t:lift}
Let $(Q,q_0)$ be pointed and $(A,B,K)$ a level-3 triple. Assume
(E) for every $F:Q\to_*K^A_3$, merely $F;K=0$; and (U) every $d:Q\to_*\Om K^B_5$
is merely $0$. Then $q_*:\hc QX\to\hc Q{K^A_3}$ is a bijection.
\end{theorem}
\begin{proof}
The text of \code{doc/mapping-fib.pdf}, \S2, with $K^A_2,K^B_4$ replaced by
$K^A_3,K^B_5$: existence from \code{Fibres.Lift}; uniqueness from (U) through
(U$_0$) ($\Om K^B_5$ connected), \code{cancel-snd}, and pointing irrelevance
($\Om X$ connected). Only connectedness of $\Om K^B_5$ and $\Om X$ is used.
Formally, \code{MapFib.Lifting} is stated at level~2, so it is either
generalised to any $k:(Y_0,y_0)\to_*(W,w_0)$ with these two connectedness
hypotheses, or copied (\S\ref{s:corr}).
\end{proof}

\section{Doubling detects the coefficient}\label{s:N}

This section is the necessity argument. Write $M_Y=\Map(K_2,Y)$, $c_2:K_2\to_*M_{K_4}$
the cup map, $\kappa=\Delta c_2$, $\iota=\cE(c_2):K_2\to_*\Om^2K_4$
(\code{cup-norm}). For $G:K_2\to_*K_2$, let $G^\sharp:M_Y\to_*M_Y$,
$m\mapsto G;m$, pointed by \code{comp-const-r}.

\begin{lemma}[the pointing of $G^\sharp$, evaluated]\label{l:ccr}
For $G:K_2\to_*K_2$ and $x:K_2$, $\ap_{\mathrm{ev}_x}(G^\sharp_0)=\refl$, where
$\mathrm{ev}_x(m)=m_1(x)$ and $G^\sharp_0=\code{comp-const-r}\,G:G;0=0$.
\end{lemma}
\begin{proof}
By definition, \code{comp-const-r} is $\ap_{\lambda e.(\lambda\_.y_0,\,e)}$ of
$\code{ap-const}:\ap_{\lambda\_.y_0}G_0=\refl$: it varies only the pointing and keeps
the underlying constant map fixed. So $\ap_{\mathrm{ev}_x}(G^\sharp_0)=
\ap_{\lambda e.\,y_0}(\code{ap-const}\ldots)$ (\code{ap-ap}), which is $\refl$ by
\code{ap-const}.
\end{proof}

\begin{lemma}[$\cE$ in the second variable]\label{l:Epre}
For $c:K_2\to_*M_Y$, $G:K_2\to_*K_2$ and $x:K_2$, $\cE(c;G^\sharp)(x)=\cE(c)(G_1x)$
in $\Om^2Y$.
\end{lemma}
\begin{proof}
The pattern of \code{ENat.e-nat} (\code{H4Bilinear}), with precomposition in place
of postcomposition; only one target $Y$, so one circle instance and no
\code{g-eq}. Write $\phi(c)=e_1^{-1};\lp c;\mathsf{ev}^{L4}:K_1\to_*N$,
$N=\Map(K_2,\Om Y)$, and $\mathrm{Ex}_x=(F\mapsto F_1x,\refl):N\to_*\Om Y$.
\emph{(1)} $\cE(c)(x)=(\lp(\phi(c);\mathrm{Ex}_x))(g)$ (\code{E-loop}: \code{ev-ap},
\code{evPt-loop}, \code{loop-comp}); likewise
$\cE(c;G^\sharp)(x)=(\lp A_1)(g)$ with $A_1=\phi(c;G^\sharp);\mathrm{Ex}_x$, and
$\cE(c)(G_1x)=(\lp A_2)(g)$ with $A_2=\phi(c);\mathrm{Ex}_{G_1x}$; both $A_i:K_1\to_*\Om Y$.
\emph{(2) An unpointed homotopy} $A_1(k)=A_2(k)$ for $k:K_1$. With
$P=(\lp c)(e_1^{-1}k)$, a loop of $M_Y$ at $0$:
$A_1(k)=\ap_{\mathrm{ev}_x}((\lp G^\sharp)P)$ (\code{ev-ap}, \code{loop-comp})
$=\ap_{\mathrm{ev}_x}\bigl((G^\sharp_0)^{-1}\cdot\ap_{G^\sharp}P\cdot G^\sharp_0\bigr)$
(\code{loop-ap}) $=\cj_{\ap_{\mathrm{ev}_x}G^\sharp_0}\bigl(\ap_{\mathrm{ev}_x\circ G^\sharp}P\bigr)$
(\code{ap-conj}, \code{ap-ap}) $=\ap_{\mathrm{ev}_{G_1x}}P$ (Lemma~\ref{l:ccr},
\code{conj-refl-loop}, and $\mathrm{ev}_x\circ G^\sharp\equiv\mathrm{ev}_{G_1x}$ by
definition of \code{ptcomp}) $=A_2(k)$ (\code{ev-ap}). This is the analogue of
\code{nat-pt}/\code{hh}; no pointing is compared.
\emph{(3) Looping the homotopy.} \code{LoopConj.LoopHomotopy} for $A_1,A_2$ and
this homotopy gives a 2-loop $\mathsf c:\Om^2Y$ with
$(\lp A_2)(g)=\cj_{\mathsf c}((\lp A_1)(g))$ in $\Om(\Om Y)$ (\code{loop-homotopy}); the conjugation is
removed by Eckmann--Hilton (\code{EH[Y,y0].eh}, \code{trans-cancel}), exactly as
in \code{e-nat}. Hence $\cE(c;G^\sharp)(x)=\cE(c)(G_1x)$.
\end{proof}

\begin{lemma}\label{l:iota2}
$\iota(K(2,2)\,y)=\iota y\cdot\iota y$ in $\Om^2K_4$, for all $y:K_2$.
\end{lemma}
\begin{proof}
$\iota$ is the inverse of $j=\lp e_3;e_2:\Om^2K_4\simeq_*K_2$
(\code{iota-eq}); apply the injective $j$. $j(\iota y\cdot\iota y)=
e_2((\lp e_3)(\iota y)\cdot(\lp e_3)(\iota y))$ (\code{loop-mult})
$=K(2,2)(e_2((\lp e_3)(\iota y)))$ (\code{Dbl.dbl} at $m=2$, $A=\Z$:
$e_2(p\cdot p)=K(2,2)(e_2p)$, by \code{sec-eq})
$=K(2,2)(j(\iota y))=K(2,2)y=j(\iota(K(2,2)y))$.
\end{proof}

\begin{proposition}[$(2x)^2=4x^2$]\label{p:Q}
$K(2,2);c_2;K(2,2)^\sharp=c_2;K(4,4)_*$ in $\Map(K_2,M_{K_4})$; hence
$K(2,2);\kappa$ and $\kappa;K(4,4)$ are pointwise equal, and
$\cls{K(2,2)};\cls\kappa=\cls\kappa;\cls{K(4,4)}$.
\end{proposition}
\begin{proof}
(\code{ETwo.e-two} and \code{Dbl.dbl} are stated with the unlocked
\code{KMap[..,twoH].K} and \code{KT.conj}; the locked \code{Lev} versions used here are
bridged by \code{K-def}/\code{conj-def}, and \code{ptcomp-assoc} puts
$K(2,2);c_2;K(2,2)^\sharp$ in the shape $K(2,2);c'$ expected by \code{e-two}; the same
bridges are used in Lemmas~\ref{l:iota2} and~\ref{l:N}.)
By Theorem~B at $B=\Z$ (\code{BilAt[Zab].theta-inj}) and \code{pi-set}
(\code{sK2}, \code{cK2}, as in \code{Kill.kill}) it suffices that the two
$\cE$ are pointwise equal. With $w=\iota y$:
$\cE(K(2,2);c_2;K(2,2)^\sharp)(y)=v\cdot v$ with $v=\cE(c_2;K(2,2)^\sharp)(y)$
(\code{ETwo.e-two}), and $v=\iota(K(2,2)y)=w\cdot w$ (Lemmas~\ref{l:Epre},
\ref{l:iota2}). On the other side, $\cE(c_2;K(4,4)_*)(y)=(\lp^2K(4,4))(w)$
(\code{ENat.e-nat}), $K(4,4)=K(2,4);K(2,4)$ (\code{k-comp}), and
$(\lp^2K(2,4))(u)=u\cdot u$ (\code{Dbl.dbl} at $m=3$ gives $\lp K(2,4)=\sq$, then
\code{loop-sq}), so it is $(w\cdot w)\cdot(w\cdot w)$. For the diagonal: by
definition $\Delta(K(2,2);c;K(2,2)^\sharp)(x)=K(2,2)^\sharp\bigl(c(K(2,2)x)\bigr)(x)
=c(K(2,2)x)(K(2,2)x)$: the precomposition already moves the second argument.
For $c=c_2$ this is $\kappa(K(2,2)x)$, and $\Delta(c_2;K(4,4)_*)(x)=K(4,4)(\kappa x)$;
evaluate the equality of bipointed maps at $x$, then at $x$. Pointings into $K_4$ are merely
irrelevant (\code{pi-mere}; $\Om K_4$ is connected), which gives the classes.
\end{proof}

\begin{lemma}[necessity]\label{l:N}
Let $\hat G:K_3\to_*K^B_5$, $H=\ucj(\lp\hat G):K_2\to_*K^B_4$, and
$h:\Hom(\Z,B)$ with $\Phi_B(h)=\cls H$. Then $h(1)+h(1)=0$.
\end{lemma}
\begin{proof}
\emph{(i) Pointwise.} $\cj H=\lp\hat G$ (\code{conj-unconj}). By
\code{Dbl.dbl}, $\sq=\cj_2K(2,2)$ on $\Om K_3$ and $\sq=\cj_4K(2_B,4)$ on
$\Om K^B_5$. So $\cj(K(2,2);H)=\cj_2K(2,2);\cj H=\sq;\lp\hat G$ and
$\cj(H;K(2_B,4))=\lp\hat G;\sq$ (associativity, $e^{-1};e=\mathrm{id}$ in the
middle). These have pointwise equal underlying maps, since
$(\lp\hat G)(p\cdot p)=(\lp\hat G)p\cdot(\lp\hat G)p$ (\code{loop-mult}).
Applying $e_4$ and $e_2^{-1}$ (\code{sec-eq}), $K(2,2);H$ and $H;K(2_B,4)$ have
pointwise equal underlying maps; by \code{pi-mere} ($\Om K^B_4$ connected) they are
merely equal, so $\cls{K(2,2)};\cls H=\cls H;\cls{K(2_B,4)}$.
\emph{(ii) Classes.} Substitute $\cls H=\cls\kappa;\cls{K(h,4)}$ (\code{Lev.K-def}
as in \code{H4Proof.phi-diag}). Left: $\cls{K(2,2)};\cls\kappa;\cls{K(h,4)}=
\cls\kappa;\cls{K(4,4)};\cls{K(h,4)}=\Phi_B(4_\Z;h)$ (Proposition~\ref{p:Q},
\code{comp-assoc}, \code{k-comp}). Right: $\cls\kappa;\cls{K(h,4)};\cls{K(2_B,4)}=\Phi_B(h;2_B)$.
By injectivity of $\Phi_B$, $h\circ4=2_B\circ h$; at $1$, with $4_\Z(1)=(1+1)+(1+1)$
(\code{twoH}, \code{hcomp}) and $h$ a homomorphism, $(h1+h1)+(h1+h1)=h1+h1$; cancelling
in $B$, $h1+h1=0$.
\end{proof}

\begin{remark}
This is exactly where a general target differs from $BR$: for $B=\Z/2$ the
conclusion is Lemma \code{two-zero}, and Proposition~\ref{p:Q} is not needed.
For $Y=\tr{S^3}_4$ the group $B=\pi_4$ is the unknown, so Lemma~\ref{l:N} cannot
be avoided. No cell of $\Unl$ is normalised, and no bipointed map is shown to
vanish from a pointwise vanishing; only classes into $K^B_4$ are compared.
\end{remark}

\begin{corollary}[the numerical converse]\label{c:conv}
For $h:\Hom(\Z,B)$ let $f_h=\cj(\kappa;K(h,4)):\Om K_3\to_*\Om K^B_5$. If
$\Unl f_h$ is inhabited, then $2h=0$. So $e_2;\kappa;K(h,4);e_4^{-1}$ deloops
if and only if $2h=0$: for ``if'', $h=\rho;g$ (Lemma~\ref{l:factor}), so
$\kappa;K(h,4)=\kb;K(g,4)$ (\code{k-comp}) and $f_h=\cj(\kb;K(g,4))=\lp(\Fh;K(g,5))$, as in
the proof of Proposition~\ref{p:EBR}.
\end{corollary}
\begin{proof}
By definition $\Unl f=\sum_{\hat G}(f=\lp\hat G)$, so an inhabitant is a pair
$(\hat G,\ \pi:f_h=\lp\hat G)$: exactly the data of Lemma~\ref{l:N}, without any
normalisation of a cell. $H=\ucj(\lp\hat G)=\ucj(f_h)=\kappa;K(h,4)$ (along
$\pi^{-1}$, and $\ucj\circ\cj=\mathrm{id}$, \code{pt-inv-l/r}, as
\code{conj-unconj}), so $\cls H=\Phi_B(h)$, and Lemma~\ref{l:N} gives
$h1+h1=0$, i.e.\ $2_B\circ h=0$ (homomorphisms out of $\Z$ are determined at $1$).
\end{proof}

\noindent What remains open from \S7 of \code{doc/cross-term.pdf} is only a
\emph{route}: normalising the cell given by \code{lemma9} at $(w_0,f_1)$, and the
general cross-term formula $\tilde\delta_F=c_2;K(2g,4)_*$. Neither is used, and
neither is needed for the numerical statement.

\begin{lemma}[factorisation through $\rho$]\label{l:factor}
If $h:\Hom(\Z,B)$ and $h1+h1=0$, then $h=\rho;g$ for $g:\Hom(\Z/2,B)$,
$g(\mathsf{false})=0$, $g(\mathsf{true})=h1$.
\end{lemma}
\begin{proof}
$g$ is a homomorphism by cases (the case $\mathsf{true}+\mathsf{true}$ uses the
hypothesis). $g(\rho x)=h(x)$ by induction on $x:\Z$ (\code{suc}, \code{pred}
as \code{rho-hom}); homomorphisms are equal when their maps are
(\code{hom-path}).
\end{proof}

\begin{proposition}[the restricted coefficient map]\label{p:Psi}
$\Psi_B:\Hom(\Z/2,B)\to\hc{K_2}{K^B_4}$, $\Psi_B(g)=\cls{\kb;K(g,4)}$, satisfies
$\Psi_B(g)=\Phi_B(\rho;g)$ and is injective. Every class $\cls{\ucj(\lp\hat G)}$
lies in its image.
\end{proposition}
\begin{proof}
$\kb;K(g,4)=\kappa;K(\rho,4);K(g,4)=\kappa;K(\rho;g,4)$ (\code{k-comp}; the same
locked \code{Lev} $K(\rho,4)$ as in $\kb$). Injective: $\rho;g=\rho;g'$ gives
$g(\mathsf{true})=g'(\mathsf{true})$ at $x=1$, and $g(\mathsf{false})=0=g'(\mathsf{false})$.
Image: Lemmas~\ref{l:N} and~\ref{l:factor}. (Conversely every
$\kb;K(g,4)$ is a looping class, of $\Fh;K(g,5)$; not needed.)
\end{proof}

\section{The source $BR$}

\begin{proposition}[$\hc R{K^B_3}=0$]\label{p:UR}
Every $d:R\to_*\Om K^B_4$ is merely $0$.
\end{proposition}
\begin{proof}
\code{doc/fibre-h3.pdf}, Theorem~10, for the triple $(\Z,\Z/2,\kb)$ in place of
$(\Z,\Z,\kappa)$: the fibre over the base point is $\Om K'_4$; the restriction
$\mathrm{res}(d)$ deloops to $G:K'_4\to_*K^B_4$ (Corollary~11, $c=2$); the
transgression gives $\cls{\kb;G}=\cls 0$; \code{Climb[Z2ab,B].surj4} gives
$G=K(g,4)$ and $0=K(g_0,4)$ with $g,g_0:\Hom(\Z/2,B)$; so
$\Psi_B(g)=\Psi_B(g_0)$, $g=g_0$ (Proposition~\ref{p:Psi}), $G=0$,
$\mathrm{res}(d)=0$, and exactness ((H3) for the base $K_2$) gives $d$ merely $0$.
The proof of \code{FibreH3} uses of $\kappa$ only the injectivity of
$g\mapsto\cls{\kappa;K(g,4)}$; it is restated with a fibre coefficient group $C$,
a map $k:K_2\to_*K^C_4$ and that injectivity as hypothesis, and instantiated at
$(C,k)=(\Z/2,\kb)$ with Proposition~\ref{p:Psi}.
\end{proof}

\begin{proposition}[(U) for $BR$]\label{p:UBR}
Every $d:BR\to_*\Om K^B_5$ is merely $0$.
\end{proposition}
\begin{proof}
$\lp:\Map(BR,\Om K^B_5)\to\Map(\Om BR,\Om^2K^B_5)$ is an equivalence
(Lemma~\ref{l:loopeq3}, $c=2$: $\Om K^B_5$ is $\cT\,6$). $\phi^{-1};\lp d:R\to_*\Om^2K^B_5$
is merely $0$ by Proposition~\ref{p:UR}, transported along
$\lp e_4:\Om^2K^B_5\simeq_*\Om K^B_4$ (\code{ptEqInd}); so merely $\lp d=0=\lp0$
(\code{pre-pteq}, \code{loop-const}), and $d=0$ by injectivity of $\lp$.
\end{proof}

\begin{lemma}[$p^*$, as \code{MapFib.SourceP.PStar}]\label{l:pstar}
The fibres of $p_1:BR\to K_3$ are $\cC\,5$ (Lemma~7 of
\code{doc/mapping-fib.pdf}, \code{fib-fst}; $K'_5$ is $\cC\,6$). For every $A$,
$p^*:\hc{K_3}{K^A_3}\to\hc{BR}{K^A_3}$ is a bijection, and every $F:BR\to_*K^A_3$
is merely $p;G_F$ with $G_F=(\mathrm{ext}\,F_1,\zeta_F)$ (\code{ConnectedMaps.Elim},
$m=5$; $K^A_3$ is $\cT\,5$; $\Om K^A_3$ connected).
\end{lemma}

\begin{lemma}[the fibre sequence]\label{l:fseq}
Merely $p;\Fh=0$ ($\lambda w.w_2$ is a homotopy; $\Om K'_5$ connected).
\end{lemma}

\begin{proposition}[(E) for $BR$]\label{p:EBR}
For every level-3 triple $(A,B,K)$ and $F:BR\to_*K^A_3$, merely $F;K=0$.
\end{proposition}
\begin{proof}
The goal is a proposition. Let $M=G_F;K:K_3\to_*K^B_5$, so merely $F;K=p;M$
(Lemma~\ref{l:pstar}, \code{ptcomp-assoc}). Let $H=\ucj(\lp M)$ and
$h=\Phi_B^{-1}\cls H$ (surjectivity of $\Phi_B$). By Lemma~\ref{l:N},
$h1+h1=0$; by Lemma~\ref{l:factor}, $h=\rho;g$; so $\cls H=\Phi_B(\rho;g)=
\cls{\kb;K(g,4)}$ (Proposition~\ref{p:Psi}) and merely $H=\kb;K(g,4)$
(\code{effective}). Then, merely,
\[\lp M=\cj H=\cj(\kb;K(g,4))=(e_2;\kb;e_4^{-1});(e_4;K(g,4);e_4^{-1})
=\lp\Fh;\lp K(g,5)=\lp(\Fh;K(g,5))\]
(\code{conj-unconj}; $e_4^{-1};e_4=\mathrm{id}$ on $K'_4$; $\gamma$; \code{loopK};
\code{loop-comp}). By Lemma~\ref{l:loopeq3} ($K_3$ is $\cC\,4$, $K^B_5$ is
$\cT\,7$), $\Unl(\lp M)$ is a proposition containing $(M,\refl)$ and
$(\Fh;K(g,5),\ldots)$; so merely $M=\Fh;K(g,5)$ (this is \code{LoopCalc.unique-loop1}
at $c=2$, which can be cited directly), and
$p;M=(p;\Fh);K(g,5)=0;K(g,5)=0$ merely (Lemma~\ref{l:fseq},
\code{comp-const-l}).
\end{proof}

\begin{corollary}\label{c:qBR}
For every level-3 triple, $q_*:\hc{BR}{\Fib K}\to\hc{BR}{K^A_3}$ is a bijection
(Theorem~\ref{t:lift}, Propositions~\ref{p:EBR}, \ref{p:UBR}).
\end{corollary}

\section{The source $S^3$ and the marked class}

\begin{lemma}\label{l:s3paths}
If $\Om^3(Z,z_0)$ is connected, any two $F,F':S^3\to_*Z$ are merely equal
($\mathrm{ev}$ is injective, Lemma~\ref{l:ev3}).
\end{lemma}

\begin{corollary}\label{c:qS}
(E) and (U) hold for $Q=S^3$ and every level-3 triple ($\Om^3K^B_5$ and
$\Om^4K^B_5$ are connected). So $q_*:\hc{S^3}{\Fib K}\to\hc{S^3}{K^A_3}$ is a
bijection; in particular $p_*:\hc{S^3}{BR}\to\hc{S^3}{K_3}$ (the triple
$(\Z,\Z/2,\Fh)$).
\end{corollary}

\begin{definition}[the marked class]\label{d:s}
$s=p_*^{-1}\cls{\sigma_3}\in\hc{S^3}{BR}$, so $s;\cls p=\cls{\sigma_3}$. A
class, obtained without choice.
\end{definition}

\begin{remark}
Under $\hc{S^3}{BR}\cong\pi_3BR\cong\pi_2R\cong\pi_2K_2\cong\Z$ the class $s$ goes
to $1$; this is not used: what is used is Proposition~\ref{p:right}, the
property that characterises the generator.
\end{remark}

\begin{proposition}\label{p:right}
For every $A$, $s^*:\hc{BR}{K^A_3}\to\hc{S^3}{K^A_3}$ is a bijection.
\end{proposition}
\begin{proof}
$s^*\circ p^*=\cls{\sigma_3}^*$ (\code{comp-assoc}); $p^*$ and
$\cls{\sigma_3}^*$ are bijections (Lemma~\ref{l:pstar},
Proposition~\ref{p:sigma3}, $K^A_3$ a 3-type); \code{cancel-r}.
\end{proof}

\section{The mapping theorem}

\begin{theorem}\label{t:main}
For every level-3 triple $(A,B,K)$, $s^*:\hc{BR}{\Fib K}\to\hc{S^3}{\Fib K}$ is
a bijection.
\end{theorem}
\begin{proof}
The square $q_*\circ s^*=s^*\circ q_*$ (\code{comp-assoc}) with top
Corollary~\ref{c:qBR}, right Proposition~\ref{p:right}, bottom
Corollary~\ref{c:qS}; \code{cancel-equiv}.
\end{proof}

\begin{corollary}[transfer]\label{c:transfer}
If merely $\Fib K\simeq_*(Y,y_0)$ for some level-3 triple, then
$s^*:\hc{BR}Y\to\hc{S^3}Y$ is a bijection (\code{ptEqInd}; the statement is a
proposition, \code{isEquiv-prop}).
\end{corollary}

\section{Postnikov presentation of 2-connected 4-types}

\begin{theorem}[\code{Presentation} one level up]\label{t:pres4}
Every pointed 2-connected 4-type $(Y,y_0)$ has merely a pointed equivalence
$\Fib K\simeq_*(Y,y_0)$ for a level-3 triple $(A,B,K)$.
\end{theorem}
\begin{proof}
The text of \code{doc/postnikov-presentation.pdf}, \S\S3--6, with indices raised,
except for the base, which is recognised instead of compared with a sphere:
\begin{enumerate}
\item $V=\tr Y_3$ (\code{TO four}, $\cT\,5$), $u:Y\to_*V$, fibres $E_z$ of $u$
  $\cC\,5$ (\code{TruncOp.conn}) and $\cT\,6$; $V$ is $\cC\,4$ (\code{conn-up}).
\item \emph{Base.} \code{Rec} (Proposition~7 there) gives $e_V:K^A_3\simeq_*V$,
  $A=\Om^3V$ (\code{O3}).
\item \emph{Fibre.} \code{Rec4}: for $F$ $\cC\,5$ and $\cT\,6$, $K^B_4\simeq_*F$ with
  $B=\Om^3(\Om F)$: \code{Rec} for $\Om F$ gives $K^B_3\simeq_*\Om F$, and one more
  application of \code{deloop-equiv} with Corollary~11 at $c=3$
  ($K^B_4$ is $\cC\,5$; $F$ is $\cT\,6$, hence $\cT\,8$), whose input is the pointed map
  $f_4=e_3;G_3:\Om K^B_4\to_*\Om F$ ($G_3$ the equivalence of \code{Rec}) together with the
  fact that $(f_4)_1$ is an equivalence (a composite). Applied to $F_0=E_{v_0}$.
\item \emph{Homogeneity.} $Q\,t=\bigl(K^B_4\simeq_*(E_{ut},(t,\refl))\bigr)$ is a set
  (\code{cor5-equiv}, $t=1$, $c=4$); a section through $Q\,t_0\ni e_{F_0}$ exists
  (\code{lemma4}, $t=0$; $Y$ is $\cC\,3$); moved along paths, $\mu_z:\mathrm{Mu}(E_z)$.
\item \emph{Classifying map.} $K^B_5$ is the type of $K^B_4$-torsors by definition of
  \code{EM} (the stage $n=2$ of $K(B,n{+}3)=T\,K(B,n{+}2)$); $k'_1(z)=\langle E_z,\mu_z\rangle$,
  $\Fib k'\simeq Y$ as Proposition~13 there; $K=e_V;k'$, $\Fib K\simeq Y$
  (\code{sigma-base}); merely pointed ($Y$ connected).
\end{enumerate}
The new pieces are \code{Rec4} and the level-2 stage of the torsor
identification (\code{K5-def}, \code{pt5-def}, $\tau$); the rest is the same text.
\end{proof}

\section{The universal property and $\pi_4(S^3)$}

\begin{theorem}[marked universal property of $BR$]\label{t:up}
For every pointed 2-connected 4-type $(Y,y_0)$, $s^*:\hc{BR}Y\to\hc{S^3}Y$ is a
bijection.
\end{theorem}
\begin{proof}
Theorem~\ref{t:pres4} and Corollary~\ref{c:transfer}.
\end{proof}

\noindent Only two instances are used. $Y=BR$: $BR=\Fib\Fh$ literally, so
Theorem~\ref{t:main} for $(\Z,\Z/2,\Fh)$ applies without Theorem~\ref{t:pres4};
in Proposition~\ref{p:EBR} the coefficient group is $\Z/2$ and Lemma~\ref{l:N} is
\code{two-zero}. $Y=T=\tr{S^3}_4$ (Proposition~\ref{p:trunc4}): Theorem~\ref{t:pres4} at
$T$, with the full Lemma~\ref{l:N}.

\begin{theorem}[$\pi_4(S^3)=\Z/2$]\label{t:final}
Under Axiom~\ref{a:s3}, $\pi_4(S^3)\cong\Z/2$ as groups.
\end{theorem}
\begin{proof}
As \code{SphereUniversal.Two}, one level up, on classes only. $u\in\hc T{BR}$
with $\cls\eta;u=s$ ($\eta^*$ at $BR$, $\cT\,6$); $t\in\hc{BR}T$ with
$s;t=\cls\eta$ (Theorem~\ref{t:up} at $T$). Then $s;(t;u)=\cls\eta;u=s=s;\mathrm{id}$,
so $t;u=\mathrm{id}_{BR}$ (Theorem~\ref{t:up} at $BR$, injective); and
$\cls\eta;(u;t)=s;t=\cls\eta=\cls\eta;\mathrm{id}$, so $u;t=\mathrm{id}_T$ ($\eta^*$ at
$T$). $\pi_4$ of a class $a$ is $\pi_3C(\lp C\,a)$, $\pi_3$ of the looped class. Obligations:
$\lp C$, looping on classes (\code{STr.Lift} of \code{Maps.loop}; the target \code{HC} is a
set); its functoriality on classes (\code{loop-comp}, \code{loop-id}); $\pi_3C$ of a class is a
homomorphism (\code{Pi3C}, already); and a level-4 interface \code{Trunc4} and \code{Two4},
replacing \code{Interfaces.Trunc3} and the hypothesis \code{isTrunc five P} of \code{Two}. So $\pi_4u$ and $\pi_4t$ are inverse isomorphisms. Then
\[\pi_4S^3\cong\pi_4T\cong\pi_4BR=\pi_3(\Om BR)\cong\pi_3R\cong\Z/2\]
(Proposition~\ref{p:trunc4}; the above; by definition; $\pi_3\phi$ for the pointed
equivalence $\phi$, \code{gi-of-equiv}; \code{Triple[Zab,Z2ab,kbar].pi3-iso}).
\end{proof}

\section{Obligations, and what to check}

\noindent\textbf{New mathematics} (to be refereed): Lemma~\ref{l:N} (necessity,
weak form), Proposition~\ref{p:Q} with Lemma~\ref{l:Epre} (second-variable
naturality of $\cE$, the one genuinely new computation; its pointing step is
flagged in the proof), Proposition~\ref{p:EBR} (the restricted lifting
argument), Proposition~\ref{p:Psi}.

\noindent\textbf{Same text, indices raised}: Lemmas~\ref{l:ev3}, \ref{l:s3conn},
Proposition~\ref{p:trunc4} (\code{SphereAxiom}, \code{SphereTruncation});
Theorem~\ref{t:lift} (\code{MapFib} \S2); Lemmas~\ref{l:pstar}, \ref{l:fseq},
Corollaries~\ref{c:qBR}, \ref{c:qS}, Theorem~\ref{t:main} (\code{MapFib}
\S\S3--5); Proposition~\ref{p:UR} (\code{FibreH3} with a fibre coefficient
group); Theorem~\ref{t:pres4} (\code{Presentation}); Theorem~\ref{t:final}
(\code{SphereUniversal.Two}).

\noindent\textbf{Settled, and not claimed.} The numerical converse of
Proposition~21 of \code{doc/cross-term.pdf} (if $\cj(\kappa;K(h,4))$ deloops then
$2h=0$) is proved here (Corollary~\ref{c:conv}). Still open, and not used: the
route through a normalised cell of \code{lemma9}, and the general cross-term
formula. Lemma~\ref{l:N} uses only pointwise information about $\lp\hat G$. No group structure on $\hc{K_2}{K^B_4}$ is used
($4h=2h$ is an equation in $\Hom(\Z,B)$). No representative of a class is
chosen.

\noindent\textbf{Points for the referee.} (1) Lemma~\ref{l:Epre}: step~(2)
uses Lemma~\ref{l:ccr} to remove the conjugation by the pointing of $G^\sharp$,
and step~(3) loops an \emph{unpointed} homotopy through \code{LoopHomotopy} and
Eckmann--Hilton, as \code{ENat.e-nat}. (2) Lemma~\ref{l:N}(i): the identification $\cj(K(2,2);H)=\sq;\cj H$
needs $\cj$ of a composite with \emph{different} groups on the two sides
($e_2$ of $\Z$, $e_4$ of $B$) and the locked \code{Lev} versions of $e_m$,
$K(2,m)$. (3) Proposition~\ref{p:EBR}: that $\Unl(\lp M)$ is a proposition
requires \code{Levels[2,1]} at target $K^B_5$ ($\cT\,7$), as in \code{DeloopR}.
(4) Theorem~\ref{t:pres4}(4): the level of \code{cor5-equiv}. (5) Whether
$\pi_4:=\pi_3\circ\Om$ agrees with the group structure intended by the final
statement (it is $\pi_1(\Om^3)$ by definition of \code{pi3}).

\section{Proposed correspondence}\label{s:corr}
\begin{itemize}
\item \code{Sphere3Axiom}: \code{ev3}, \code{Sphere3[S3,base,surf,ump3]}:
  \code{evPt}, \code{pointed-up}, \code{evPt-nat} (Lemma~\ref{l:ev3}).
\item \code{Sphere3Trunc}: \code{s3-conn} (Lemma~\ref{l:s3conn}), \code{Trunc4}
  record and \code{trunc4} (Proposition~\ref{p:trunc4}); \code{pi4}, \code{pi4-map},
  \code{loopC}, \code{pi4C} in \code{HomotopyGroups} style.
\item \code{Necessity}: \code{ccr-ev} (Lemma~\ref{l:ccr}), \code{E2}, \code{A1}, \code{A2}, \code{hh}, \code{e-pre} (Lemma~\ref{l:Epre}), \code{iota-dbl}
  (Lemma~\ref{l:iota2}), \code{quad}, \code{kappa-dbl} (Proposition~\ref{p:Q}),
  \code{N.pointwise}, \code{N.cls}, \code{N.two-tors} (Lemma~\ref{l:N}), \code{conv} (Corollary~\ref{c:conv}),
  \code{factor} (Lemma~\ref{l:factor}), \code{Psi}, \code{psi-phi}, \code{psi-inj}
  (Proposition~\ref{p:Psi}).
\item \code{FibreH3}: \code{FibAtC[C,k,inj,B]} (the existing proof with a fibre
  coefficient group); \code{fib-h3-R} (Proposition~\ref{p:UR}).
\item \code{MapBR}: \code{Triple3[A,B,K]} (fibre $\Fib K$, its pointing, \code{conn4} via
  \code{FibConn}, \code{trunc6}, connectedness of $\Om\Fib K$), \code{Conj[eA,eB]}
  (Conventions), \code{Lifting3} (Theorem~\ref{t:lift}); \code{SourceBR}:
  \code{cp}, \code{PStar}, \code{fseq}, \code{ExBR}, \code{ex-BR}
  (Proposition~\ref{p:EBR}), \code{un-BR} (Proposition~\ref{p:UBR}), \code{qBR};
  \code{SourceS3}: \code{s3-paths}, \code{qS}, \code{sigma3}, \code{Sig3}
  (Definition~\ref{d:sigma3}, Proposition~\ref{p:sigma3}); \code{Final}: \code{s},
  \code{s-eq}, \code{right}, \code{rep-fib}, \code{rep-pteq}.
\item \code{Presentation4}: \code{Rec4}, \code{Pres4[Y,y0,cY,tY]}, \code{pres4}
  (Theorem~\ref{t:pres4}).
\item \code{Pi4Sphere}: \code{up} (Theorem~\ref{t:up}), \code{Two4},
  \code{pi4-S3 : GIso[pi4 S3 base, ab-grp Z2ab]} (Theorem~\ref{t:final}).
\end{itemize}
Each declaration gets an entry in a new spec file (\code{br.opens}), every line
``nothing else opened''.
\end{document}
